\subsection*{Components (H4, H5)}
Table~\ref{tab:abl} (and Fig.~\ref{fig:abl}, right) varies one component at a time at $K\le 100$ with the tuned configuration, five seeds. The reference row (``Conditional AR designer'') repeats the main configuration with a maximum budget of $K=100$, so its $K=1$ value differs slightly from Table~\ref{tab:main} through sampling noise only (same trained model, different random draws).

\begin{table*}[t]\centering\caption{Ablations of the AR designer (five seeds). test\_nll: held-out NLL of the true designing sequences (nats/residue).}\label{tab:abl}
\begin{tabular}{lccccccccccccccccc}
\toprule
Method & succ\_k1 $\uparrow$ & succ\_k10 $\uparrow$ & succ\_k100 $\uparrow$ & gap\_k1 $\downarrow$ & gap\_k10 $\downarrow$ & gap\_k100 $\downarrow$ & succ\_k10\_single $\uparrow$ & succ\_k10\_multi $\uparrow$ & succ\_k100\_single $\uparrow$ & succ\_k100\_multi $\uparrow$ & test\_nll $\downarrow$ & sample\_rate $\uparrow$ & train\_pairs $\uparrow$ & n\_targets $\uparrow$ & n\_train\_conf $\uparrow$ & n\_single\_targets $\uparrow$ & $n$ \\
\midrule
\textbf{Conditional AR designer} & \textbf{0.60 $\pm$ 0.07} & \underline{0.80 $\pm$ 0.04} & \underline{0.87 $\pm$ 0.03} & \textbf{0.01 $\pm$ 0.01} & \underline{0.00 $\pm$ 0.00} & \textbf{0.00 $\pm$ 0.00} & \textbf{0.75 $\pm$ 0.08} & \underline{0.83 $\pm$ 0.05} & \underline{0.82 $\pm$ 0.06} & \underline{0.90 $\pm$ 0.04} & \underline{0.33 $\pm$ 0.05} & \textbf{0.60 $\pm$ 0.04} & \textbf{1954.80 $\pm$ 115.94} & \textbf{90.00 $\pm$ 0.71} & \textbf{276.00 $\pm$ 0.71} & \textbf{28.80 $\pm$ 4.32} & 5 \\
No reversal augmentation & \underline{0.53 $\pm$ 0.03} & \textbf{0.81 $\pm$ 0.03} & \textbf{0.91 $\pm$ 0.01} & \underline{0.02 $\pm$ 0.01} & \textbf{0.00 $\pm$ 0.00} & \underline{0.00 $\pm$ 0.00} & \underline{0.74 $\pm$ 0.07} & \textbf{0.85 $\pm$ 0.05} & \textbf{0.86 $\pm$ 0.05} & \textbf{0.93 $\pm$ 0.02} & \textbf{0.27 $\pm$ 0.03} & \underline{0.55 $\pm$ 0.04} & 977.40 $\pm$ 57.97 & \underline{90.00 $\pm$ 0.71} & \underline{276.00 $\pm$ 0.71} & \underline{28.80 $\pm$ 4.32} & 5 \\
Unconditional AR & 0.00 $\pm$ 0.00 & 0.00 $\pm$ 0.00 & 0.01 $\pm$ 0.00 & 3.77 $\pm$ 0.28 & 1.86 $\pm$ 0.17 & 0.69 $\pm$ 0.09 & 0.00 $\pm$ 0.00 & 0.00 $\pm$ 0.00 & 0.00 $\pm$ 0.00 & 0.01 $\pm$ 0.01 & 0.73 $\pm$ 0.03 & 0.00 $\pm$ 0.00 & \underline{1954.80 $\pm$ 115.94} & 90.00 $\pm$ 0.71 & 276.00 $\pm$ 0.71 & 28.80 $\pm$ 4.32 & 5 \\
\bottomrule
\end{tabular}

\end{table*}

\textbf{H4 (conditioning): supported.} The unconditional AR model reaches \rhval{abl_components/unconditional-ar/hp16/succ_k100/mean:3} at $K=100$ versus \rhval{abl_components/conditional-ar-designer/hp16/succ_k100/mean:3} for the conditional model (paired $p=$ \rhval{cmp/abl_components/unconditional-ar/hp16/succ_k100/paired_p}), i.e.\ a generative model of designable sequences alone is no better than random search at this budget (Table~\ref{tab:main}); the information in the target contacts is what makes the designs work. We did not test which feature group (contact rows, free-neighbour counts) is responsible.

\textbf{H5a (reversal augmentation).} The registered comparison (above, full data) is uninformative about augmentation. Splits keep each chain-reversal class together, so at data fraction 1.0 the reverse of every training pair is already in the training set: ``augmentation'' duplicates the pairs (about \rhval{analysis_aug/augmentation-analysis/hp16/f1p0_aug_train_pairs/mean:0} pairs, of which \rhval{analysis_aug/augmentation-analysis/hp16/f1p0_aug_distinct_pairs/mean:0} distinct, on average over seeds; the distinct count was not logged for the original runs, but equals the pair count of the unaugmented model, and an audit confirmed that all augmented pairs are duplicates). The ``No reversal augmentation'' row therefore compares identical data trained for twice as many gradient steps with the ``augmented'' (duplicated) set, i.e.\ it measures training length. With that caveat, the registered numbers are: $K=1$ success \rhval{abl_components/conditional-ar-designer/hp16/succ_k1/mean:3} with against \rhval{abl_components/no-reversal-augmentation/hp16/succ_k1/mean:3} without (difference with minus without \rhval{cmp/abl_components/no-reversal-augmentation/hp16/succ_k1/delta:3}, paired $p=$ \rhval{cmp/abl_components/no-reversal-augmentation/hp16/succ_k1/paired_p:3}); at $K=100$ the model without \emph{duplication} is better (\rhval{abl_components/no-reversal-augmentation/hp16/succ_k100/mean:3} against \rhval{abl_components/conditional-ar-designer/hp16/succ_k100/mean:3}, difference \rhval{cmp/abl_components/no-reversal-augmentation/hp16/succ_k100/delta:3}, $p=$ \rhval{cmp/abl_components/no-reversal-augmentation/hp16/succ_k100/paired_p:3}, uncorrected), with a lower held-out NLL (\rhval{abl_components/no-reversal-augmentation/hp16/test_nll/mean:3} against \rhval{abl_components/conditional-ar-designer/hp16/test_nll/mean:3}). A step-matched control (no augmentation, 120 epochs) on the same data gives \rhval{analysis_aug/augmentation-analysis/hp16/f1p0_noaug120_succ_k1_mean/mean:3} at $K=1$ and \rhval{analysis_aug/augmentation-analysis/hp16/f1p0_noaug120_succ_k100_mean/mean:3} at $K=100$, indistinguishable from the ``augmented'' model (Table~\ref{tab:aug}, last row), which confirms that the registered difference is a training-length effect. Epochs were not tuned for the two settings.

\textbf{H5a, follow-up where augmentation adds data (post hoc).} To test augmentation proper we re-ran the ablation at data fractions 0.1, 0.25 and 0.5, where a random subset of the sequence space is scored and only part of each reversal pair survives (Table~\ref{tab:aug}); these runs were added after the audit of the registered result and are exploratory. Against the unaugmented model at equal epochs (control A), augmentation raises $K=1$ success at fractions 0.1 and 0.25 (differences \rhval{analysis_aug/augmentation-analysis/hp16/f0p1_noaug60_succ_k1_delta/mean:3}, $p=$ \rhval{analysis_aug/augmentation-analysis/hp16/f0p1_noaug60_succ_k1_p/mean:3}; and \rhval{analysis_aug/augmentation-analysis/hp16/f0p25_noaug60_succ_k1_delta/mean:3}, $p=$ \rhval{analysis_aug/augmentation-analysis/hp16/f0p25_noaug60_succ_k1_p/mean:3}; uncorrected, five seeds), not detectably at 0.5 ($p=$ \rhval{analysis_aug/augmentation-analysis/hp16/f0p5_noaug60_succ_k1_p/mean:3}), and shows no detectable change in $K=100$ success at any of the three (Table~\ref{tab:aug}). Against the step-matched control B, augmented success is higher at $K=100$ for all three fractions (Table~\ref{tab:aug}), most at fraction 0.1. Control B shows that simply training the unaugmented model longer on its smaller set does not recover the augmented model's quality; it does not separate the new pairs from the extra gradient signal they bring, and we did not test other regularisers. We read the evidence as: augmentation helps in the data-poor regime, mostly at $K=1$ at equal epochs, its benefit shrinks as the data grows, and at full data it is a no-op apart from the longer training. Even the follow-up is a single architecture with untuned epochs.

\begin{table*}[t]\centering\caption{Reversal augmentation across data fractions (five seeds; $\Delta$ = augmented minus control, paired $t$-test $p$ in parentheses, uncorrected). ``Distinct'' counts unique (features, sequence) pairs in the augmented set. Control A: no augmentation, 60 epochs (half the gradient steps of the augmented model). Control B: no augmentation, 120 epochs (same number of gradient steps). At fraction 1.0 augmentation adds no new pair, so the comparison with control B is the same data and the same steps.}\label{tab:aug}
\begin{tabular}{lrrrrrrr}\hline
Fraction & Distinct & Pairs aug. & Pairs none & $\Delta$K1 vs A & $\Delta$K100 vs A & $\Delta$K1 vs B & $\Delta$K100 vs B \\\hline
0.1 & \rhval{analysis_aug/augmentation-analysis/hp16/f0p1_aug_distinct_pairs/mean:0} & \rhval{analysis_aug/augmentation-analysis/hp16/f0p1_aug_train_pairs/mean:0} & \rhval{analysis_aug/augmentation-analysis/hp16/f0p1_noaug_train_pairs/mean:0} & \rhval{analysis_aug/augmentation-analysis/hp16/f0p1_noaug60_succ_k1_delta/mean:3} (\rhval{analysis_aug/augmentation-analysis/hp16/f0p1_noaug60_succ_k1_p/mean:3}) & \rhval{analysis_aug/augmentation-analysis/hp16/f0p1_noaug60_succ_k100_delta/mean:3} (\rhval{analysis_aug/augmentation-analysis/hp16/f0p1_noaug60_succ_k100_p/mean:3}) & \rhval{analysis_aug/augmentation-analysis/hp16/f0p1_noaug120_succ_k1_delta/mean:3} (\rhval{analysis_aug/augmentation-analysis/hp16/f0p1_noaug120_succ_k1_p/mean:3}) & \rhval{analysis_aug/augmentation-analysis/hp16/f0p1_noaug120_succ_k100_delta/mean:3} (\rhval{analysis_aug/augmentation-analysis/hp16/f0p1_noaug120_succ_k100_p/mean:3}) \\
0.25 & \rhval{analysis_aug/augmentation-analysis/hp16/f0p25_aug_distinct_pairs/mean:0} & \rhval{analysis_aug/augmentation-analysis/hp16/f0p25_aug_train_pairs/mean:0} & \rhval{analysis_aug/augmentation-analysis/hp16/f0p25_noaug_train_pairs/mean:0} & \rhval{analysis_aug/augmentation-analysis/hp16/f0p25_noaug60_succ_k1_delta/mean:3} (\rhval{analysis_aug/augmentation-analysis/hp16/f0p25_noaug60_succ_k1_p/mean:3}) & \rhval{analysis_aug/augmentation-analysis/hp16/f0p25_noaug60_succ_k100_delta/mean:3} (\rhval{analysis_aug/augmentation-analysis/hp16/f0p25_noaug60_succ_k100_p/mean:3}) & \rhval{analysis_aug/augmentation-analysis/hp16/f0p25_noaug120_succ_k1_delta/mean:3} (\rhval{analysis_aug/augmentation-analysis/hp16/f0p25_noaug120_succ_k1_p/mean:3}) & \rhval{analysis_aug/augmentation-analysis/hp16/f0p25_noaug120_succ_k100_delta/mean:3} (\rhval{analysis_aug/augmentation-analysis/hp16/f0p25_noaug120_succ_k100_p/mean:3}) \\
0.5 & \rhval{analysis_aug/augmentation-analysis/hp16/f0p5_aug_distinct_pairs/mean:0} & \rhval{analysis_aug/augmentation-analysis/hp16/f0p5_aug_train_pairs/mean:0} & \rhval{analysis_aug/augmentation-analysis/hp16/f0p5_noaug_train_pairs/mean:0} & \rhval{analysis_aug/augmentation-analysis/hp16/f0p5_noaug60_succ_k1_delta/mean:3} (\rhval{analysis_aug/augmentation-analysis/hp16/f0p5_noaug60_succ_k1_p/mean:3}) & \rhval{analysis_aug/augmentation-analysis/hp16/f0p5_noaug60_succ_k100_delta/mean:3} (\rhval{analysis_aug/augmentation-analysis/hp16/f0p5_noaug60_succ_k100_p/mean:3}) & \rhval{analysis_aug/augmentation-analysis/hp16/f0p5_noaug120_succ_k1_delta/mean:3} (\rhval{analysis_aug/augmentation-analysis/hp16/f0p5_noaug120_succ_k1_p/mean:3}) & \rhval{analysis_aug/augmentation-analysis/hp16/f0p5_noaug120_succ_k100_delta/mean:3} (\rhval{analysis_aug/augmentation-analysis/hp16/f0p5_noaug120_succ_k100_p/mean:3}) \\
1.0 & \rhval{analysis_aug/augmentation-analysis/hp16/f1p0_aug_distinct_pairs/mean:0} & \rhval{analysis_aug/augmentation-analysis/hp16/f1p0_aug_train_pairs/mean:0} & \rhval{analysis_aug/augmentation-analysis/hp16/f1p0_noaug_train_pairs/mean:0} & \rhval{analysis_aug/augmentation-analysis/hp16/f1p0_noaug60_succ_k1_delta/mean:3} (\rhval{analysis_aug/augmentation-analysis/hp16/f1p0_noaug60_succ_k1_p/mean:3}) & \rhval{analysis_aug/augmentation-analysis/hp16/f1p0_noaug60_succ_k100_delta/mean:3} (\rhval{analysis_aug/augmentation-analysis/hp16/f1p0_noaug60_succ_k100_p/mean:3}) & \rhval{analysis_aug/augmentation-analysis/hp16/f1p0_noaug120_succ_k1_delta/mean:3} (\rhval{analysis_aug/augmentation-analysis/hp16/f1p0_noaug120_succ_k1_p/mean:3}) & \rhval{analysis_aug/augmentation-analysis/hp16/f1p0_noaug120_succ_k100_delta/mean:3} (\rhval{analysis_aug/augmentation-analysis/hp16/f1p0_noaug120_succ_k100_p/mean:3}) \\
\hline\end{tabular}\end{table*}

\subsection*{Data size (H5b)}
We score only a random fraction of the sequence space and train on the uniquely folding sequences found, which are the training pairs (Fig.~\ref{fig:abl}, left). Averaged over five seeds, $K=1$ success rises with the fraction at every step in the means: \rhval{analysis/data-fraction-analysis/hp16/frac0p1_succ_k1_mean/mean:3} (fraction 0.1, about \rhval{analysis/data-fraction-analysis/hp16/frac0p1_train_pairs_mean/mean:0} pairs after augmentation, about \rhval{analysis_aug/augmentation-analysis/hp16/f0p1_aug_distinct_pairs/mean:0} distinct), \rhval{analysis/data-fraction-analysis/hp16/frac0p25_succ_k1_mean/mean:3} (0.25), \rhval{analysis/data-fraction-analysis/hp16/frac0p5_succ_k1_mean/mean:3} (0.5) and \rhval{analysis/data-fraction-analysis/hp16/frac1p0_succ_k1_mean/mean:3} (1.0, about \rhval{analysis/data-fraction-analysis/hp16/frac1p0_train_pairs_mean/mean:0} pairs after augmentation but only \rhval{analysis_aug/augmentation-analysis/hp16/f1p0_aug_distinct_pairs/mean:0} distinct). Success at $K=100$ is nearly flat: \rhval{analysis/data-fraction-analysis/hp16/frac0p1_succ_k100_mean/mean:3}, \rhval{analysis/data-fraction-analysis/hp16/frac0p25_succ_k100_mean/mean:3}, \rhval{analysis/data-fraction-analysis/hp16/frac0p5_succ_k100_mean/mean:3}, \rhval{analysis/data-fraction-analysis/hp16/frac1p0_succ_k100_mean/mean:3}. The verdict is mixed: more scored sequences improve the per-sample quality of the model at $K=1$ (the change from fraction 0.1 to 1.0 is large against the seed std of about \rhval{analysis/data-fraction-analysis/hp16/frac1p0_succ_k1_std/mean:3}; we ran no paired tests on the individual steps between neighbouring fractions, which are smaller), but with an oracle that can filter 100 samples the amount of data barely matters here. A confound: at fixed epochs, fewer pairs also means fewer gradient steps, and Table~\ref{tab:aug} shows that the number of steps matters, so the sweep does not separate data from training length. We did not test data fractions above 1.0 because the space is exhausted.

\begin{figure*}[t]\centering\includegraphics[width=0.85\textwidth]{ablations.pdf}
\caption{Left: AR success versus the fraction of the sequence space scored (training-pair count grows with it), at $K=1$ and $K=100$. Right: component ablations. Mean $\pm$ std over five seeds.}\label{fig:abl}
\end{figure*}

\subsection*{What the ablations do not show}
The sweep and ablations use the hyperparameters tuned for the full model; the unconditional model in particular was not tuned separately, although its failure ($\approx$ random level) is too large to be a tuning effect. The seeds control both the split and the model initialisation, so seed variance mixes target-set difficulty with training noise; we did not separate them.
